\chapter{KIỂM THỬ HỒI QUY VÀ ĐÁNH GIÁ CHẤT LƯỢNG}

\section{Thực hiện Kiểm thử hồi quy (Regression Testing)}

\subsection{Định nghĩa và Tầm quan trọng}
Kiểm thử hồi quy là hoạt động kiểm thử lại toàn bộ hoặc một phần hệ thống sau khi có sự thay đổi mã nguồn (như sửa lỗi hoặc thêm tính năng mới) nhằm đảm bảo rằng các sửa đổi này không làm phát sinh lỗi mới tại các mô-đun chức năng vốn đã chạy ổn định trước đó.

\subsection{Quy trình vá lỗi của Lập trình viên và kiểm thử lại của Tester}
Sau khi nhận được 3 ticket báo lỗi nghiêm trọng từ Tester trên Mantis Bug Tracker, Lập trình viên đã tiến hành vá lỗi trực tiếp trên mã nguồn 7coMART:
\begin{enumerate}
    \item \textbf{Vá lỗi bảo mật SQL Injection (Ticket \#0000021):} Lập trình viên đã thay thế truy vấn nối chuỗi SQL thô bằng việc sử dụng ORM SQLAlchemy, tự động tham số hóa (Parameterized Query) mọi biến đầu vào của người dùng, triệt tiêu hoàn toàn khả năng tiêm nhiễm mã SQL.
    \item \textbf{Vá lỗi tìm kiếm tiếng Việt (Ticket \#0000022):} Bổ sung mã hóa ký tự UTF-8 cho kết nối Database trong file cấu hình \texttt{core/database.py} của FastAPI và cấu hình lại kiểu dữ liệu đối chiếu (Collation) của trường tên sản phẩm trong PostgreSQL thành \texttt{"Vietnamese\_Vietnam.1258"} (hoặc UTF-8 mặc định).
    \item \textbf{Vá lỗi số lượng âm giỏ hàng (Ticket \#0000023):} Sử dụng thư viện Pydantic trong Python FastAPI để bổ sung điều kiện ràng buộc dữ liệu đầu vào (\texttt{Field(gt=0)}) tại Schema của CartItem, đảm bảo số lượng gửi lên bắt buộc phải lớn hơn 0 ở mức API.
\end{enumerate}
Tester tiến hành lấy bản build Docker mới nhất đã sửa đổi, chạy lại toàn bộ bộ kịch bản test (15 Test Cases). Kết quả xác nhận cả 3 lỗi nghiêm trọng đều đã được sửa đúng, các tính năng đăng nhập, tìm kiếm, giỏ hàng hoạt động hoàn toàn chính xác. Tester tiến hành chuyển trạng thái 3 ticket lỗi trên Mantis sang \textbf{Closed}.

\section{Phân tích các độ đo chất lượng phần mềm nâng cao}
Để đánh giá hiệu quả của quy trình kiểm thử và chất lượng phần mềm website 7coMART một cách khách quan, nhóm áp dụng hai độ đo toán học tiêu chuẩn công nghiệp:

\subsection{Độ đo Mật độ lỗi (Defect Density - DD)}
Mật độ lỗi thể hiện số lượng lỗi phát hiện được trên mỗi mô-đun chức năng kiểm thử của hệ thống, giúp định vị các khu vực kém ổn định để tập trung nguồn lực tối ưu hóa. Công thức tính:
$$DD = \frac{N_{\text{bugs}}}{N_{\text{modules}}}$$
\textit{Trong đó:}
\begin{itemize}
    \item $N_{\text{bugs}}$: Tổng số lỗi phát hiện được trong quá trình kiểm thử (Nhóm ghi nhận tổng số 8 lỗi).
    \item $N_{\text{modules}}$: Tổng số mô-đun chức năng tham gia kiểm thử (Đăng ký, Đăng nhập, Tìm kiếm, Giỏ hàng - gồm 4 mô-đun lớn).
\end{itemize}
\textit{Tính toán thực tế:}
$$DD = \frac{8}{4} = 2.0 \text{ (lỗi / mô-đun)}$$
Chỉ số mật độ lỗi bằng 2.0 phản ánh hệ thống ở mức trung bình, các lỗi nghiêm trọng đã được kiểm soát tốt sau pha hồi quy.

\subsection{Độ đo Hiệu quả loại bỏ lỗi (Defect Removal Efficiency - DRE)}
DRE là chỉ số đo lường tỷ lệ lỗi được phát hiện và loại bỏ trong các giai đoạn kiểm thử nội bộ so với tổng số lỗi thực tế tồn tại trong sản phẩm (bao gồm cả các lỗi rò rỉ bị người dùng phát hiện sau khi bàn giao). Công thức:
$$DRE = \frac{E}{E + D} \times 100\%$$
\textit{Trong đó:}
\begin{itemize}
    \item $E$: Số lượng lỗi được phát hiện và xử lý triệt để trong pha kiểm thử nội bộ của nhóm Tester (8 lỗi).
    \item $D$: Số lượng lỗi rò rỉ phát sinh do người dùng phát hiện ở pha vận hành thực tế (Giả định mô phỏng ghi nhận 1 lỗi nhỏ về giao diện bị bỏ sót).
\end{itemize}
\textit{Tính toán thực tế:}
$$DRE = \frac{8}{8 + 1} \times 100\% = 88.89\%$$
Chỉ số DRE đạt $88.89\%$, vượt qua ngưỡng tiêu chuẩn chất lượng của dự án đồ án tốt nghiệp ($>85\%$), minh chứng cho năng lực thiết kế test case và kiểm thử hộp đen xuất sắc của nhóm sinh viên.

\section{Báo cáo kiểm thử tổng hợp (Test Summary Report)}
Dưới đây là bảng tổng hợp kết quả thực thi kiểm thử 7coMART, phân loại chi tiết theo từng mô-đun và thống kê số lượng lỗi theo các mức độ nghiêm trọng khác nhau.

\begin{table}[H]
\centering
\caption{Bảng tổng kết thực thi kiểm thử theo từng mô-đun chức năng}
\renewcommand{\arraystretch}{1.2}
\begin{tabular}{|l|c|c|c|c|}
\hline
\textbf{Mô-đun kiểm thử} & \textbf{Tổng số TC} & \textbf{Số TC Pass} & \textbf{Số TC Fail} & \textbf{Tỷ lệ Đạt (\%)} \\ \hline
Đăng ký tài khoản & 4 & 4 & 0 & 100.0\% \\ \hline
Đăng nhập hệ thống & 3 & 2 & 1 & 66.7\% \\ \hline
Tìm kiếm \& Lọc Catalog & 3 & 2 & 1 & 66.7\% \\ \hline
Giỏ hàng \& Đặt hàng & 5 & 4 & 1 & 80.0\% \\ \hline
\textbf{Tổng cộng} & \textbf{15} & \textbf{12} & \textbf{3} & \textbf{80.0\%} \\ \hline
\end{tabular}
\end{table}

\begin{table}[H]
\centering
\caption{Thống kê số lượng lỗi phát hiện phân loại theo Mức độ nghiêm trọng}
\renewcommand{\arraystretch}{1.2}
\begin{tabular}{|l|c|c|c|}
\hline
\textbf{Mức độ nghiêm trọng} & \textbf{Số lỗi phát hiện} & \textbf{Đã khắc phục (Closed)} & \textbf{Tỷ lệ vá lỗi (\%)} \\ \hline
Block (Ngăn chặn bảo mật) & 1 & 1 & 100.0\% \\ \hline
Critical (Lỗi nghiệp vụ nặng) & 1 & 1 & 100.0\% \\ \hline
Major (Lỗi chức năng lớn) & 2 & 2 & 100.0\% \\ \hline
Minor/Tweak (Lỗi nhỏ/Giao diện) & 4 & 4 & 100.0\% \\ \hline
\textbf{Tổng cộng} & \textbf{8} & \textbf{8} & \textbf{100.0\%} \\ \hline
\end{tabular}
\end{table}

% =========================================================================
% HƯỚNG DẪN CHÈN ẢNH CHO SINH VIÊN:
% Xuất báo cáo đồ thị (Summary chart) phân loại lỗi hình bánh/cột trong MantisBT,
% lưu ảnh thành file 'mantis_report_chart.png' và bỏ vào figs/
% =========================================================================
\begin{figure}[H]
    \centering
    \begin{tikzpicture}[scale=0.8, every node/.style={transform shape}]
        % Browser frame
        \draw[thick, fill=white] (0,0) rectangle (15,9.5);
        \draw[fill=gray!20, thick] (0,8.7) rectangle (15,9.5);
        \fill[red] (0.3,9.1) circle (0.12);
        \fill[yellow] (0.7,9.1) circle (0.12);
        \fill[green] (1.1,9.1) circle (0.12);
        \draw[fill=white, rounded corners=2pt] (2.0,8.85) rectangle (12.0,9.35);
        \node[right] at (2.2,9.1) {\scriptsize \texttt{http://localhost:8089/summary\_page.php}};
        
        % Left Sidebar (MantisBT menu)
        \draw[fill=gray!80, draw=none] (0,0) rectangle (2.5,8.7);
        \node[white, right] at (0.3,8.0) {\tiny $\bullet$ Xem báo cáo thống kê};
        
        % Main Workspace Header
        \node[right] at (2.8,8.2) {\small\textbf{Thống kê và phân tích báo cáo sự cố (7coMART)}};
        \draw[gray!30] (2.8,7.9) -- (14.5,7.9);
        
        % Pie Chart container
        \draw[draw=gray, fill=white] (3.0,1.0) rectangle (8.5,7.2);
        \node at (5.75,6.7) {\tiny\textbf{Tỷ lệ Lỗi theo Mức độ Nghiêm trọng}};
        
        % Drawing Pie Chart
        \begin{scope}[shift={(5.75,3.6)}]
            % Block (12.5%): 0 to 45 deg
            \draw[fill=red!60, draw=white, thick] (0,0) -- (0:2) arc (0:45:2) -- cycle;
            % Critical (12.5%): 45 to 90 deg
            \draw[fill=orange!70, draw=white, thick] (0,0) -- (45:2) arc (45:90:2) -- cycle;
            % Major (25%): 90 to 180 deg
            \draw[fill=yellow!70, draw=white, thick] (0,0) -- (90:2) arc (90:180:2) -- cycle;
            % Minor/Tweak (50%): 180 to 360 deg
            \draw[fill=blue!40, draw=white, thick] (0,0) -- (180:2) arc (180:360:2) -- cycle;
        \end{scope}
        
        % Pie Chart Legend
        \fill[red!60] (3.5,2.5) rectangle (3.8,2.8); \node[right] at (3.9,2.65) {\tiny Block (1 lỗi - 12.5\%)};
        \fill[orange!70] (3.5,2.0) rectangle (3.8,2.3); \node[right] at (3.9,2.15) {\tiny Critical (1 lỗi - 12.5\%)};
        \fill[yellow!70] (3.5,1.5) rectangle (3.8,1.8); \node[right] at (3.9,1.65) {\tiny Major (2 lỗi - 25.0\%)};
        \fill[blue!40] (3.5,1.0) rectangle (3.8,1.3); \node[right] at (3.9,1.15) {\tiny Minor/Tweak (4 lỗi - 50.0\%)};
        
        % Statistics Summary Table
        \draw[draw=gray, fill=gray!5] (9.0,1.0) rectangle (14.0,7.2);
        \node at (11.5,6.7) {\tiny\textbf{Bảng Tổng Hợp Số Liệu}};
        \draw[gray!30] (9.2,6.3) -- (13.8,6.3);
        
        \node[right] at (9.2,5.9) {\tiny $\bullet$ Tổng số sự cố: 8};
        \node[right] at (9.2,5.4) {\tiny $\bullet$ Đã xử lý (Closed): 8 (100\%)};
        \node[right] at (9.2,4.9) {\tiny $\bullet$ Mật độ lỗi: 2.0 lỗi/mô-đun};
        \node[right] at (9.2,4.4) {\tiny $\bullet$ Hiệu quả loại bỏ lỗi (DRE): 88.89\%};
        
        \node[right] at (9.2,3.4) {\tiny $\bullet$ Số lượng Tester: 4};
        \node[right] at (9.2,2.9) {\tiny $\bullet$ Số lượng Developer: 1};
        \node[right] at (9.2,2.4) {\tiny $\bullet$ Sprints hoàn thành: 2};
        
        \draw[draw=green, fill=green!10, rounded corners=3pt] (9.2,1.2) rectangle (13.8,2.0);
        \node[green!60!black, right] at (9.5,1.6) {\tiny\textbf{Trạng thái chất lượng:}};
        \node[green!60!black, right] at (9.5,1.3) {\tiny\textbf{ĐẠT YÊU CẦU BÀN GIAO}};
    \end{tikzpicture}
    \caption{Biểu đồ thống kê tỷ lệ phần trăm phân bố trạng thái lỗi và mức độ nghiêm trọng trích xuất từ hệ thống Mantis Bug Tracker}
    \label{fig:mantis_report_chart}
\end{figure}

\textbf{Nhận xét kết quả:} Như minh họa tại biểu đồ Hình \ref{fig:mantis_report_chart}, tỷ lệ lỗi nghiêm trọng (Block, Critical) chỉ chiếm góc nhỏ khoảng 25\% tổng số lượng lỗi, trong khi các lỗi vừa và nhỏ (Major, Minor) chiếm 75\%. Việc phân tích cấu trúc lỗi này giúp đội ngũ quản lý chất lượng khẳng định ứng dụng 7coMART đã đạt độ ổn định rất cao sau pha hồi quy, đủ điều kiện để đóng dự án và triển khai thực tế.